\subsection{Producing a homotopy inverse from collapsible intersections}
\label{sec:collapse-nerve}

The linear contraction in Lemma~\ref{lem:contraction-producer} fills
chains. A geometric contraction allows one to fill maps, which gives a
stronger conclusion for a decidable class of covers. Use the subcomplex
cover $K=\bigcup_i K_i$, intersections $K_I$, and nerve $N$ of
Section~\ref{sec:nerve}. A legal elementary collapse removes a pair
$(\sigma,\tau)$, where $\sigma$ is nonempty, $\tau=\sigma\cup\{v\}$
for one new vertex $v$, and $\tau$ is the only face strictly containing
$\sigma$ in the current complex. Call a complex collapsible if a finite
sequence of such removals leaves a single vertex.

\begin{theorem}\label{thm:collapse-nerve}
There is a primitive-recursive test whether every nonempty $K_I$ is
collapsible. On success it produces rational evaluators with PR moduli
for maps $|K|\rightleftarrows |N|$ and homotopies from both composites to
their respective identities. On failure it identifies an intersection
without a collapse sequence of the indicated type.
\end{theorem}
\begin{proof}
\emph{Local production.}
If an intersection has $S$ nonempty faces, each move removes two.
Enumerate words of length at most $S$ over at most $S^2$ possible
pairs, check each move in its current face table, and test the final
vertex condition. The number of words is at most $(1+S^2)^{S+1}$.
This search includes every successful collapse sequence. Perform it on
each of the at most $2^r-1$ nonempty index sets, omitting empty
intersections. An unsuccessful search excludes collapsibility for this
triangulation; a different triangulation or a different contraction
can have different certificates.

One elementary move supplies an actual deformation to the remaining
subcomplex. For barycentric coordinates $x$ in the current complex put
$m(x)=\min_{j\in\sigma}x_j$. At time $t$, subtract $t m(x)$ from each
coordinate in $\sigma$ and add $|\sigma|t m(x)$ to the coordinate $v$.
If $m(x)>0$, the support contains $\sigma$, so it lies in $\tau$.
The altered point stays in that simplex, with nonnegative coordinates
of sum one. At $t=1$ at least one coordinate in $\sigma$ vanishes. If
the input lies in the remaining subcomplex its support misses a vertex
of $\sigma$, so the entire motion fixes it. These assertions also cover
the case where $\sigma$ is a vertex. In the sup norm on both coordinates
and time, the move has Lipschitz bound $2|\sigma|+1$: the minimum changes
by at most the coordinate error, and multiplication by $t\in[0,1]$
adds at most the time error.

If a sequence has length $s$, apply its moves in order at respective
times $t_j=\min(1,\max(0,st-j))$, for $0\le j<s$.
All preceding moves are complete whenever a move is active. Inactive
moves have coefficient zero in the same coordinate formula. This
constructs a rational contraction $H_I(x,t)$ from identity to its
terminal vertex. For input on at most $v_0$ vertices, a common
Lipschitz bound in the product sup norm is obtained by $B_0=1$,
$B_{j+1}=(2v_0+1)(B_j+s+1)$, taking $B_s$ at the end. A zero-length
sequence occurs on a single vertex and uses the identity there.

\emph{Extension of maps in nested carriers.}
Here a carrier for a simplex $e$ of a finite complex $L$ is a
subcomplex $C(e)$ of the target with $C(e')\subseteq C(e)$ for every
nonempty face $e'$ of $e$. Suppose each $C(e)$ comes with a rational
contraction $H_e$ from identity to a specified point and with a Lipschitz
bound. A map already defined on a subcomplex, with restrictions Lipschitz
and carried simplexwise, extends constructively over $L$. We describe
the construction because it supplies the uniformity needed in iPRA.

Initialize missing vertices at the contraction endpoints. For a missing
simplex with $n\ge2$ vertices the entire boundary has already been mapped
into $C(e)$. Write its map as $b$ and suppose its restrictions to facets
have Lipschitz bound $L_0\ge1$. The whole boundary has bound $2nL_0$
in barycentric sup norm. Indeed, if rational points $x,y$ have distance
$\delta$ with $n\delta\ge1$, all target coordinates lie in $[0,1]$.
Otherwise their supports meet: the mass outside the common support is at
most $n\delta$. Form $z$ by taking coordinatewise minima and moving
the remaining mass, at most $n\delta$, to any common vertex.
Both $x,z$ lie in one source facet and $y,z$ in another;
$\|x-z\|_\infty,\|y-z\|_\infty\le n\delta$.
Agreement of the facet maps on their common face yields the bound.
The same argument glues simplexwise Lipschitz bounds over any finite
source complex, using its total number of vertices.

Write $c=(1/n,\ldots,1/n)$ and
$u(x)=\max_i(1-n x_i)$. It ranges from $0$ at $c$ to $1$ on the
boundary. For $u>1/2$ put $y=c+(x-c)/u$; the choice of $u$ ensures
that the coordinates of $y$ are nonnegative, sum to one, and include
a zero. Extend by
\[
  E(x)=
  \begin{cases}
    H_e(b(y),2(1-u)),&u>1/2,\\
    H_e(\text{endpoint},1),&u\le1/2.
  \end{cases}
\]
This agrees with $b$ at $u=1$ and is constant on the inner half-sized
simplex. For two outer inputs of distance $\delta$, their radial
parameters differ by at most $n\delta$, their inverses by at most
$4n\delta$, and the boundary arguments differ by at most
$(2+4n)\delta$. At an inner/outer pair, compare the outer expression
with $H_e(b(y),1)$ using the same outer boundary argument; the time
difference is at most $2n\delta$. Thus, if $B$ bounds the contraction,
\[
 B\bigl[2nL_0(2+4n)+2n+1\bigr]
\]
bounds the extension in all cases. At a rational boundary argument one
may evaluate any least indexed facet containing it, since the facet
maps agree there. The threshold comparison takes place on rational
samples. The displayed modulus extends the result to names. Visiting
simplices in increasing dimension proves the extension assertion.

The same construction gives a carried homotopy between two carried maps,
with known Lipschitz bounds, by extending their prescribed values at
the two ends of a cylinder. To make the cylinder data explicit, order
the vertices of $L$ once. For a simplex on $x_0<\cdots<x_d$ triangulate
its cylinder with the simplices
\[
 (x_0,0),\ldots,(x_j,0),(x_j,1),\ldots,(x_d,1)
 \qquad(0\le j\le d)
\]
and their faces. The assignment restricts compatibly on faces. The
carrier of a cylinder simplex is $C(e)$ for its projected source face
$e$. The identification of the product with this triangulation splits
the ordered barycentric masses at total mass $1-t$: assign the initial
mass to time-zero vertices and the remainder to time-one vertices,
using the $b_i$ formula in Theorem~\ref{thm:nerve-pl}. Its support
is contained in one displayed simplex. The inverse adds the two
coordinates at each source vertex and takes the total top mass as time.
Both operations have explicit Lipschitz bounds by finite summation and
minimum. Thus the extension produces a homotopy with the stated
endpoints also across every shared face.

\emph{The two comparison maps.}
For a face $\sigma$ of $K$, put $F(\sigma)=\{i:\sigma\in K_i\}$.
This is a nerve face and $F$ reverses inclusion, giving a simplicial
map $F:|\operatorname{sd}K|\to|\operatorname{sd}N|$ by accumulating
weights. Identify subdivisions with their originals by the rational
coordinate formulas of Theorem~\ref{thm:nerve-pl}. To construct
$G:|\operatorname{sd}N|\to |K|$, give a chain
$e=(I_0\subsetneq\cdots\subsetneq I_d)$ the carrier $K_{I_0}$.
Removing its initial vertex can only shrink the carrier. Apply the
extension just proved, starting with its chosen vertices. In particular
$G(e)\subseteq |K_{I_0}|$, with $G(e)$ here denoting the image of
the whole geometric simplex.

For the composite $GF$, a source chain
$\sigma_0\subsetneq\cdots\subsetneq\sigma_d$ has carrier
$K_{F(\sigma_d)}$. Both its identity image inside $\sigma_d$ and the
image under $GF$ lie there: all $F(\sigma_j)$ contain $F(\sigma_d)$,
and the image of the entire chain is a face of the reverse chain with
least element $F(\sigma_d)$. These carriers are nested on faces and
have the produced contractions. The cylinder construction gives the
first homotopy.

For $FG$ use the order complex of the poset of all nerve faces
containing $I_0$ as carrier for $e$. It contracts to the vertex $I_0$
by $(1-t)x+t\mathbf e_{I_0}$: adjoining the least element to any chain
still gives a chain. Its Lipschitz bound is at most three. The identity
is carried there. Every point of $|K_{I_0}|$, when expressed in
subdivision coordinates, has its positively weighted faces $\sigma$
in $K_{I_0}$. For rational points this follows from the explicit
coordinate formula: each such $\sigma$ is a subset of the point's
support. Thus $F(\sigma)\supseteq I_0$, and $FG$ is carried there too.
Apply the cylinder construction again. All verifications on rational
points pass to names with the prescribed moduli and the resulting
uniform equality estimates.

Finally, one may bound the number of subdivision vertices by
$D=\max(1,2^{v_0},2^r)$ and the number of simplices needed in any
one cylinder by $2^{2D}$. At each inductive stage take the maximum
of the preceding bound and the displayed extension bound over the
finite list, then proceed to the next dimension. A rational point
evaluation follows at most one descending face chain per extension;
each step uses a bounded facet scan, rational arithmetic and a
contraction with a previously bounded number of moves. The successive
constructions of $G$, the endpoint composites and the two homotopies
have fixed nesting depth as constructions. Table enumeration, the
dimension recurrence for the Lipschitz constants and these evaluation
clocks are PR. An output tolerance $2^{-k}$ is obtained by evaluating
a rational input sample to tolerance $2^{-k}/(2L)$ for the resulting
bound $L\ge1$. Empty realizations use the empty maps. This proves
the claim for the represented spaces.
\end{proof}

In the hexagon example each two-edge arc collapses by removing its
two outer vertices with their incident edges, leaving its middle vertex.
Each nonempty pairwise intersection already consists of one vertex.
The theorem therefore upgrades the computed homology transfer for that
cover to a homotopy equivalence with the triangular nerve.

For Lean, add to the chain-comparison targets three independent
finite kernels: the free-face test and min-coordinate collapse, the
collar extension with its two cases, and the cylinder incidence list.
Record the extension invariant as agreement on previously processed
faces, membership in the assigned carrier, and the displayed Lipschitz
estimate. For an object proof, keep the two input sample points and the
requested accuracy as parameters. For each fixed construction, carry the
invariant through the descending evaluator path and its bounded table
lookups; a finite trace records all rationals and comparisons used in
both evaluations. Equivalently, bound their codes in advance by iterating
a majorant for one rational operation through the PR operation clock,
then Booleanize the invariant for all such subcalls at each depth.
The geometric maps and homotopies in this section are proved on paper;
the supplied \texttt{ElementaryCollapse.lean} checks mass preservation, nonnegativity, the free-face endpoint, and the fixed-wall property of the local deformation. The carrier extensions and their assembly remain deep formalization targets.
